% Einheit 1 von 5 · Klammern auflösen (terme.md, Einheit 3) – Hauptnummern 8–15
\setcounter{aufgabe}{7}
\zweigkopf{Einheit 1 von 5 · Klammern auflösen}{Hier lernst du, Klammern mit Plus, Minus oder einer Zahl davor aufzulösen · neu oder schon bekannt – je nach Buch · keine P10-Aufgabe · baut auf: Terme zusammenfassen, Malnehmen mit Vorzeichen}{e1}

\begin{aufgabe}{Ich erkenne die Glieder eines Terms mit ihrem Vorzeichen. Schreibe die Glieder mit ihrem Vorzeichen auf.}
\begin{geruest}
\gz{a}{$5x - 2y + 7$}{\leerfeld\leerfeld\leerfeld}
\gz{b}{$-3m + 4 - n$}{\leerfeld\leerfeld\leerfeld}
\gz{c}{$x^2 - 6x - 1$}{\leerfeld\leerfeld\leerfeld}
\gz{d}{$-8 + 2x - x^2$}{\leerfeld\leerfeld\leerfeld}
\end{geruest}
\end{aufgabe}

\begin{aufgabe}{Ich erkenne, was vor einer Klammer steht. Schreibe auf, was direkt vor der Klammer steht: Plus, Minus, eine positive Zahl oder eine negative Zahl. Löse nichts auf.}
\begin{geruest}
\gz{a}{$9 + (2x - 5)$}{\leerfeld}
\gz{b}{$6 - (x + 4)$}{\leerfeld}
\gz{c}{$3 \cdot (2x - 1)$}{\leerfeld}
\gz{d}{$8x - 4 \cdot (x - 2)$}{\leerfeld}
\gz{e}{$2x - (5 - x)$}{\leerfeld}
\end{geruest}
\end{aufgabe}

\begin{aufgabe}{Ich kann eine Klammer mit Plus davor weglassen. Löse die Klammer auf und fasse zusammen.}
\beispiel{6 + (x - 1) &= 6 + x - 1 \\ &= x + 5}
\begin{gleichungsraster}[2]
\gl{3 + (x + 4)} & \gl{2x + (5 + x)} \\ \noalign{\vspace{5mm}}
\gl{10 + (x - 6)} & \gl{4m + (2 - m)} \\ \noalign{\vspace{5mm}}
\end{gleichungsraster}
\end{aufgabe}

\begin{aufgabe}{Ich kann eine Zahl mit einer Klammer malnehmen. Multipliziere jedes Glied der Klammer mit der Zahl.}
\beispiel{4 \cdot (x + 2) &= 4x + 8}
\par\smallskip Als Tabelle: \quad\sachtabelle{ccc}{$\cdot$ & $x$ & $2$}{$4$ & $4x$ & $8$}
\begin{gleichungsraster}[1]
\gl{2 \cdot (x + 3)} & \gl{5 \cdot (m + 4)} \\ \noalign{\vspace{5mm}}
\gl{6 \cdot (y + 1)} & \gl{7 \cdot (2 + x)} \\ \noalign{\vspace{5mm}}
\gl{4 \cdot (x - 5)} & \gl{(x + 9) \cdot 2} \\ \noalign{\vspace{5mm}}
\gl{x \cdot (x + 4)} & \\ \noalign{\vspace{5mm}}
\end{gleichungsraster}
\end{aufgabe}

\begin{aufgabe}{Ich kann eine Klammer mit Minus davor auflösen (neu in diesem Jahr). Löse die Klammer auf; vor der Klammer steht ein Minus, also drehst du jedes Vorzeichen in der Klammer. Fasse dann zusammen.}
\beispiel{9 - (x - 2) &= 9 - x + 2 \\ &= 11 - x}
\begin{gleichungsraster}[2]
\gl[1]{-(x + 6)} & \gl{10 - (x + 2)} \\ \noalign{\vspace{5mm}}
\gl{7 - (x - 5)} & \gl{5x - (2x - 1)} \\ \noalign{\vspace{5mm}}
\gl{4m - (m - 2n + 3)} & \gl{-3 \cdot (x - 4)} \\ \noalign{\vspace{5mm}}
\gl{4x - 2 \cdot (x + 3)} & \gl[3]{2 \cdot (3x + 5) - 4 \cdot (x - 2)} \\ \noalign{\vspace{5mm}}
\end{gleichungsraster}
\end{aufgabe}

\begin{aufgabe}{Ich finde den Fehler beim Auflösen von Klammern. Finde den Fehler, benenne ihn und korrigiere die Rechnung. Rechne danach die Aufgabe darunter selbst.}
\begin{teile}
\teil Jonas rechnet so:
\end{teile}
\rechnung{12 - (x - 5) &= 12 - x - 5 \\ &= 7 - x}
\begin{gleichungsraster}[2]
\gl{9 - (2x - 4)} & \\ \noalign{\vspace{5mm}}
\end{gleichungsraster}
\begin{teile}
\teil Nina rechnet so:
\end{teile}
\rechnung{5 - 2 \cdot (x + 1) &= 3 \cdot (x + 1) \\ &= 3x + 3}
\begin{gleichungsraster}[2]
\gl{8 - 3 \cdot (x + 2)} & \\ \noalign{\vspace{5mm}}
\end{gleichungsraster}
\end{aufgabe}

\begin{aufgabe}{Ich kann begründen, ob zwei Terme gleichwertig sind. Gleichwertig heißt: Für jede Zahl, die man für $x$ einsetzt, haben beide Terme denselben Wert.}
\begin{teile}
\teil Sind $3 \cdot (x - 2)$ und $3x - 2$ gleichwertig? Begründe.
\teil Sind $6 - (x + 2)$ und $4 - x$ gleichwertig? Begründe.
\teil Lea setzt $x = 1$ in $2 \cdot (x + 1)$ und in $x + 3$ ein und erhält beide Male $4$. Sie sagt: „Die Terme sind gleichwertig.“ Hat sie recht? Begründe.
\end{teile}
\end{aufgabe}

\begin{aufgabe}{Ich kann eine Sachaufgabe mit einer Klammer lösen. Lisa kauft für jedes ihrer drei Geschwister ein Heft zu $p$~€ und einen Stift zu $2$~€. Sie bezahlt mit einem 20-€-Schein.}
\begin{teile}
\teil Stelle einen Term für das Wechselgeld auf.
\teil Löse die Klammer auf und fasse zusammen.
\teil Ein Heft kostet $1{,}50$~€. Wie viel Wechselgeld bekommt Lisa?
\teil Reicht das Geld, wenn ein Heft $5$~€ kostet? Begründe mit dem Term.
\end{teile}
\end{aufgabe}
